\section{Signed moments in the complete count record}\label{sec:count}

The realization has more structure than its selected no-count output.
For each $j\ge0$, define
\begin{equation}\label{eq:countdef}
 \sigma_j(t)=K_j(t,0)\rho_*K_j(t,0)^*,\qquad
 P_j(t)=\Tr\sigma_j(t),\qquad
 d_j(t,k)=\Tr[\sigma_j(t)\cos(kQ_t)].
\end{equation}
Thus $P_0=p_t$ and $d_0$ is the absolute contrast in \eqref{eq:heat}.
For nonzero counts these are well-defined analytical traces of the existing
output operators. The frozen physical contract supplies the selected
no-count readout; mathematical definition of $d_j$ does not add a guarantee
of access to the same measurement after every nonzero count.

The literal count integral and normalization below were specifically
checked in \cite{TNDHighCountReview}. Section~\ref{sec:highcount} gives an
independent assessed complete-source proof of count summation and its first
real-$k$ derivative \cite{TNDCountTailReview}. The finite transfer laws and
whole-channel kernel calculation below retain their proof-author
diagnostic status; these limited assessments do not cover every claim of
\cite{TNDDirectCollision}.

\begin{theorem}[Count-sector moment and generating identities]\label{thm:counts}
For fixed $j\ge0$, $0\le t\le T$ and real $k$,
\begin{align}
 P_j(t)&=p_t\frac{(\gamma t)^j}{j!}\mu_{2j}(t),\label{eq:Pj}\\
 d_j(t,k)&=\frac{(\gamma t)^j}{j!Z_{1/2}}
      \int_\R v^{2j}F_{\lambda(t)}(v)\cos(kv)\dd v
       =\frac{(-\gamma t)^j}{j!}\partial_k^{2j}d_0(t,k),\label{eq:dj}\\
 |d_j(t,k)|&\le P_j(t),\qquad\sum_{j\ge0}P_j(t)=1.\label{eq:countabs}
\end{align}
At $t=0$ the stated count-zero convention applies. For $0\le\eta\le1$,
\begin{equation}\label{eq:generate}
 \sum_{j\ge0}\eta^jd_j(t,k)
 =\frac{H_{\lambda(t)+\eta\gamma t}(k)}{H_{1/2}(0)}.
\end{equation}
The series converges absolutely. In particular,
\begin{equation}\label{eq:conserved}
 \sum_{j\ge0}d_j(t,k)=C(0,k).
\end{equation}
\end{theorem}
\begin{proof}
In the interaction frame the radial output vector is
$\ee^{-\gamma tv^2/2}(\sqrt{\gamma t}\,v)^jf_{1/2}(v)/\sqrt{j!}$.
The signed odd-count amplitude is kept; its two occurrences in the diagonal
trace give $v^{2j}$. The complete angular trace is exactly one. This proves
the first integral formulas, including \eqref{eq:Pj}. Taking $2j$ derivatives
in $k$ is justified by \eqref{eq:Hmixed} and produces $(-1)^jv^{2j}$,
proving \eqref{eq:dj}. Positivity of $\sigma_j$ and the contraction norm of
the real cosine give \eqref{eq:countabs}; count completeness gives the sum.

Inside the source integral, the absolute sum of the terms weighted by
$\eta^j$ is dominated by the positive sum with the cosine replaced by one.
It is at most the initial probability mass for $0\le\eta\le1$.
The exponential series therefore gives
\[
 \frac1{Z_{1/2}}\int_\R F_{\lambda(t)}(v)
        \ee^{\eta\gamma tv^2}\cos(kv)\dd v
 =\frac{2H_{\lambda(t)+\eta\gamma t}(k)}{Z_{1/2}},
\]
which is \eqref{eq:generate}. At $\eta=1$ the parameter returns to $1/2$,
proving \eqref{eq:conserved}. No derivative-order summation is inferred
from fixed-order estimates: this series is justified by its own absolute
probability bound.
\end{proof}

Equation~\eqref{eq:Pj} is also a mixture of Poisson count laws against the
initial density, with scalar intensity $\gamma tv^2$. For $t>0$, every
$P_j>0$ and its normalized signed trace is
\begin{equation}\label{eq:conditionalj}
 \frac{d_j}{P_j}=\frac{(-1)^j\partial_k^{2j}C}{\mu_{2j}}.
\end{equation}
The count record therefore organizes the same complete weighted moments
that direct source differentiation supplies.

\begin{proposition}[Finite-sector transfer laws]\label{prop:transfer}
For fixed $j$ and $t>0$,
\begin{equation}\label{eq:transfer}
 \partial_td_j=\frac jt d_j-\frac{j+1}{t}d_{j+1},\qquad
 \partial_tP_j=\frac jt P_j-\frac{j+1}{t}P_{j+1}.
\end{equation}
Let $q_{j,\pm}=(P_j\pm d_j)/2$. These nonnegative analytical effect
probabilities obey the same equations, and in particular
\begin{equation}\label{eq:loss}
 \partial_tq_{0,\pm}=-q_{1,\pm}/t\le0,\qquad
 \partial_td_0=-d_1/t.
\end{equation}
\end{proposition}
\begin{proof}
Differentiate the integral in \eqref{eq:dj}. The prefactor contributes
$jd_j/t$; differentiating $F_{\lambda(t)}$ contributes $-\gamma v^2$,
whose integral is $(j+1)d_{j+1}/t$ with the stated sign. The same calculation
without the cosine proves the formula for $P_j$. Finite moments justify
each fixed derivative. Taking sum and difference gives the equations for
$q_{j,\pm}$. No infinite-series derivative is used.
\end{proof}

The quotients in \eqref{eq:transfer} have the required limits at zero from
the source formulas; they should not be evaluated as $0/0$. For example,
$d_1(t,k)/t\to\gamma\E_0[X^2\cos(kX)]$. Every statement concerns the
fixed slab and fixed positive $\gamma$, without an infinite-time or uniform
small-$\gamma$ conclusion.

\subsection{What is conserved when all counts are retained}
The conservation in \eqref{eq:conserved} concerns the co-moving cosine
observation. It does not make the full state stationary. The distinction
can be seen directly from the same Kraus amplitudes. In the interaction
frame, the all-count output of the prescribed initial state has the radial
operator-valued kernel
\begin{equation}\label{eq:dephase}
 f_{1/2}(v)f_{1/2}(w)\,\ee^{-\gamma t(v-w)^2/2}R.
\end{equation}
Indeed summing the kernel multipliers gives
\[
 \ee^{-\gamma t(v^2+w^2)/2}
       \sum_{j\ge0}\frac{(\gamma tvw)^j}{j!}
 =\ee^{-\gamma t(v-w)^2/2}.
\]
The absolute series is bounded by
$\exp[-\gamma t(|v|-|w|)^2/2]\le1$, so the kernels converge in
Hilbert--Schmidt norm by domination by $|f_{1/2}(v)f_{1/2}(w)|$.
Theorem~\ref{thm:instrument} already gives trace-norm convergence to the
same operator. The diagonal density is unchanged, while off-diagonal
coherence is attenuated. This kernel identity is an algebraic consequence
of the fixed instrument, not an additional physical evolution assumption.
It also explains why replacing signed odd-count amplitudes by their
absolute values would change the complete operation.
